\documentclass[11pt]{article}
% \usepackage[top=0.4in, bottom=0.2in, left=0.75in, right=0.75in]{geometry}
\usepackage{amssymb}
\usepackage{amsthm}
\usepackage[pdftex]{graphicx}
\usepackage{epsfig}
\usepackage{amstext}
\usepackage{amsmath}
\usepackage{multicol}
\usepackage[T1]{fontenc}
\usepackage{yfonts}
\usepackage[numbers,sort]{natbib}
\bibliographystyle{abbrvnat}
\usepackage{natbib}
\usepackage{hyperref}
\usepackage{listings}
\usepackage{xcolor}
\usepackage{algorithm,algpseudocode}

\usepackage{subcaption} % Required for the subfigure environment
\usepackage{float}      % Required for the [H] positioning specifier
\oddsidemargin -0.2in \textwidth 7in \topmargin -1in \textheight 9.25in

\newtheorem{theorem}{Theorem}[section]
\newtheorem{lemma}[theorem]{Lemma}
\newtheorem{corollary}[theorem]{Corollary}
\newtheorem{remark}[theorem]{Remark}
\newtheorem{definition}[theorem]{Definition}
\newtheorem{assumption}[theorem]{Assumption}

% \title{\textbf{Data-Driven Kernel Reconstruction and Threshold Dynamics}}
% \author{
%     Jose Aries E. De Los Santos \\
%     \textit{jedelossantos1@up.edu.ph}
% }
% \date{April 21, 2026}
% \begin{document}
% \maketitle
% \thispagestyle{empty}
% --- MANUAL HEADER (Replacing \maketitle) ---

\begin{document}
\begin{center}
{\LARGE \textbf{A Modified Parameterized Gaussian Error Linear Unit} \\[1em]
{\large \textbf{Jose Aries E. De Los Santos}$^{1,2}$ \quad | \quad \textit{jedelossantos1@up.edu.ph}} \\[0.5em]
{\normalsize 
$^1$Institute of Mathematics, College of Science, University of the Philippines Diliman \\
$^2$Data Science Program, College of Science, University of the Philippines Diliman}
}
\end{center}
\hrule \vspace{0.5em}

\section{Background}
\begin{definition}
The normal distribution with parameter values $\mu = 0$ and $\sigma = 1$ is called the standard normal distribution. A random variable having a standard normal distribution is called a standard normal random variable and will be denoted by $Z$. The probability density function of $Z$ is
\begin{equation*}
f(z) = \frac{1}{\sqrt{2\pi}}e^{-z^2/2}
\end{equation*}
where $z \in \mathbb{R}$. The cumulative distribution function of $Z$ is
$$
\mathbb{P}(Z \leq z) = \int_{-\infty}^{z} f(t)dt
$$
which we denote by $\Phi(z)$.
\end{definition}
\subsection{Gaussian Error Linear Unit (GELU) Formulation}
\noindent We first observe the similarities between the Rectified Linear Unit (ReLU)~\cite{nair2010rectified} and the Heaviside function.
\begin{equation*}
\text{ReLU}(x) = \max(0, x) = 
\begin{cases} 
    x, & \text{if } x > 0 \\ 
    0, & \text{if } x \leq 0 
\end{cases}
\qquad \text{and} \qquad
H(x) = 
\begin{cases} 
    1, & \text{if } x > 0 \\ 
    0, & \text{if } x \leq 0 
\end{cases}
\end{equation*}
\noindent As seen in their definitions, the Rectified Linear Unit (ReLU) can be explicitly expressed as the input $x$ multiplied by the Heaviside step function $H(x)$, yielding $\text{ReLU}(x) = x \cdot H(x)$. This acts as a hard, deterministic gating mechanism. In contrast, regularization techniques like Dropout apply a stochastic gate, scaling the input by a random mask $m \sim \text{Bernoulli}(p)$. The stochastic dropout operation $\nu(x)$ is formulated as:
$$
\nu(x) = m \cdot x = 
\begin{cases} 
    x, & \text{with probability } p \\ 
    0, & \text{with probability } 1 - p 
\end{cases}
$$
The expected value of this stochastic transformation is $\mathbb{E}[\nu(x)] = x \cdot p + 0 \cdot (1-p) = x \cdot p$. Motivated by this, Hendrycks and Gimpel~\cite{hendrycks2023gaussianerrorlinearunits} sought to merge these paradigms by applying a zero-one mask that is stochastically determined yet strictly dependent upon the input. They defined this stochastic activation as:
$$
f(x) = m \cdot x
$$
where $m \sim \text{Bernoulli}(\Phi(x))$. Rather than dropping out neurons uniformly with a constant probability $p$, the Gaussian Error Linear Unit (GELU) defines the probability of preserving the input as a function of the input itself, specifically utilizing the standard normal cumulative distribution function. By replacing $p$ with $\Phi(x)$, GELU is mathematically formulated as the expected value of this input-dependent stochastic process: 
$$
\text{GELU}(x) = \mathbb{E}[m \cdot x] = x \Phi(x).
$$
Here, $\Phi(x)$ is the cumulative distribution function of the standard normal distribution, defined as:
$$
\Phi(x) = \mathbb{P}(X \leq x) = \frac{1}{\sqrt{2\pi}} \int_{-\infty}^{x} e^{-t^2/2} dt.
$$
This specific distribution is used because neuron inputs naturally tend to follow a normal distribution, particularly when the network applies Batch Normalization. To evaluate the cumulative distribution function, we can split $\Phi(x)$ into two distinct intervals at zero:
\begin{equation*}
\Phi(x) = \frac{1}{\sqrt{2\pi}} \int_{-\infty}^{x} e^{-t^2/2} dt =  \frac{1}{\sqrt{2\pi}} \int_{-\infty}^{0} e^{-t^2/2} dt + \frac{1}{\sqrt{2\pi}} \int_{0}^{x} e^{-t^2/2} dt
\end{equation*}
The first term represents the left half of the probability distribution. To understand its exact value, we must first determine the total area under the foundational Gaussian curve:
\begin{equation*}
\int_{-\infty}^{\infty} e^{-x^2} \, dx
\end{equation*}
We start by squaring the integral and introducing two independent dummy variables, $x_1$ and $x_2$. Fubini's Theorem allows us to combine the product of two one-dimensional integrals into a single double integral over the entire Cartesian plane:
\begin{align*}
\left( \int_{-\infty}^{\infty} e^{-x^2} \, dx \right)^2 &= \left( \int_{-\infty}^{\infty} e^{-x_1^2} \, dx_1 \right) \left( \int_{-\infty}^{\infty} e^{-x_2^2} \, dx_2 \right) \\
&= \int_{-\infty}^{\infty} \int_{-\infty}^{\infty} e^{-(x_1^2 + x_2^2)} \, dx_1 \, dx_2
\end{align*}
Next, we apply a change of variables to polar coordinates. Recall that for a continuous function on a polar region, 
\begin{equation*}
\iint_R f(x,y) dA = \int_{\alpha}^{\beta} \int_a^b f(r\cos\theta, r\sin\theta) r \, dr \, d\theta
\end{equation*}
Applying the substitutions $x_1^2 + x_2^2 = r^2$ and $dx_1 \, dx_2 = r \, dr \, d\theta$, and expanding the bounds to cover the entire infinite plane ($r$ from $0$ to $\infty$ and $\theta$ from $0$ to $2\pi$), we obtain:
\begin{equation*}
\int_{-\infty}^{\infty} \int_{-\infty}^{\infty} e^{-(x_1^2 + x_2^2)} \, dx_1 \, dx_2 = \int_0^{2\pi} \int_0^{\infty} e^{-r^2} r \, dr \, d\theta
\end{equation*}
We evaluate the inner integral with respect to $r$ using a standard $u$-substitution (where $u = -r^2$ and $du = -2r \, dr$):
\begin{align*}
\int_0^{2\pi} \int_0^{\infty} e^{-r^2} r \, dr \, d\theta &= \int_0^{2\pi} \left[ -\frac{1}{2} e^{-r^2} \right]_{r=0}^{r=\infty} \, d\theta \\
&= \int_0^{2\pi} \left[ 0 - \left(-\frac{1}{2}\right) \right] \, d\theta \\
&= \int_0^{2\pi} \frac{1}{2} \, d\theta
\end{align*}
Evaluating the final outer integral with respect to $\theta$:
\begin{equation*}
\int_0^{2\pi} \frac{1}{2} \, d\theta = \left[ \frac{1}{2}\theta \right]_0^{2\pi} = \pi
\end{equation*}
Thus, the square of our original integral equals $\pi$. Taking the square root of both sides yields the foundational result:
\begin{equation*}
\left( \int_{-\infty}^{\infty} e^{-x^2} \, dx \right)^2 = \pi \implies \int_{-\infty}^{\infty} e^{-x^2} \, dx = \sqrt{\pi}
\end{equation*}
Now, we apply this result to find the total area under our unscaled standard normal curve:
\begin{equation*}
\text{Total Area} = \int_{-\infty}^{\infty} e^{-z^2/2} \, dz
\end{equation*}
We evaluate this using another $u$-substitution. Let $x = \frac{z}{\sqrt{2}}$, which implies $z = x\sqrt{2}$ and $dz = \sqrt{2} \, dx$. Substituting this into the integral and keeping the infinite bounds, we get:
\begin{equation*}
\int_{-\infty}^{\infty} e^{-x^2} (\sqrt{2} \, dx) = \sqrt{2} \int_{-\infty}^{\infty} e^{-x^2} \, dx
\end{equation*}
For a function to qualify as a valid probability density function, the total area under the curve must integrate to exactly $1$. Dividing the raw kernel by its total area establishes our normalization factor:
\begin{equation*}
\int_{-\infty}^{\infty} \dfrac{1}{\sqrt{2\pi}} e^{-z^2/2} \, dz = \dfrac{1}{\sqrt{2\pi}} \cdot \sqrt{2\pi} = 1
\end{equation*}
Because the curve is perfectly symmetric around zero and the total normalized area is exactly $1$, the integral of the left half (from $-\infty$ to $0$) evaluates to precisely to:
\begin{equation*}
\dfrac{1}{\sqrt{2\pi}} \int_{-\infty}^{0} e^{-z^2/2} \, dz = \dfrac{1}{2}
\end{equation*}
This allows us to resolve the first term of our split $\Phi(x)$ function, simplifying it to:
\begin{equation*}
\Phi(x) = \frac{1}{2} + \frac{1}{\sqrt{2\pi}} \int_{0}^{x} e^{-t^2/2} dt
\end{equation*}
While the integral for $\Phi(x)$ perfectly describes the theory, practical implementations in deep learning frameworks (such as PyTorch~\cite{paszke2017automatic}) rely on the Error Function:
\begin{equation*}
\text{erf}(z) = \dfrac{2}{\sqrt{\pi}}\int_{0}^{z} e^{-u^{2}}du
\end{equation*}
Let $u = \frac{t}{\sqrt{2}}$. Then, $t = u\sqrt{2}$ and $dt = \sqrt{2} \, du$. 
For the limits of integration, we have:
\begin{equation*}
\begin{aligned}
\text{If } t &= 0 \text{ then } u = 0 \\
\text{If } t &= x \text{ then } u = \frac{x}{\sqrt{2}}
\end{aligned}
\end{equation*}
Substituting this into $\frac{1}{\sqrt{2\pi}} \int_{0}^{x} e^{-t^2/2} dt$, we obtain:
\begin{equation}
\frac{1}{\sqrt{2\pi}} \int_{0}^{x} e^{-t^2/2} dt 
= \frac{1}{\sqrt{2\pi}} \int_{0}^{x/\sqrt{2}} e^{-u^2} (\sqrt{2} \, du) 
= \frac{1}{\sqrt{\pi}} \int_{0}^{x/\sqrt{2}} e^{-u^{2}} du 
\label{gauss_int}
\end{equation}
And we can express~\eqref{gauss_int} as:
\begin{equation*}
\frac{1}{\sqrt{\pi}} \int_{0}^{x/\sqrt{2}} e^{-u^{2}} du = \frac{1}{2} \left[ \frac{2}{\sqrt{\pi}} \int_{0}^{x/\sqrt{2}} e^{-u^{2}} du \right] = \frac{1}{2}\text{erf}\left(\frac{x}{\sqrt{2}}\right) 
\end{equation*}
Finally, we obtain:
\begin{equation*}
\text{GELU}(x) = x\Phi(x) =  \frac{x}{2}\left[1 + \text{erf}\left( \frac{x}{\sqrt{2}}\right) \right]
\end{equation*}
\noindent The error function is a standard mathematical tool specifically built to calculate areas under Gaussian curves, and it is highly optimized at the hardware level. By rewriting our cumulative distribution function $\Phi(x)$ in terms of the error function, we can compute the exact GELU efficiently in software.
\section{MPGELU Formulation}

\noindent The standard Gaussian Error Linear Unit (GELU) relies on the fixed variance of the standard normal distribution $\mathcal{N}(0, 1)$, which limits the network's ability to adapt to varying data distributions across layers. To address this, we investigate a parameterized extension that introduces a learnable scaling parameter $\lambda \geq 1$. This formulation independently aligns with the recently proposed $\lambda$-GELU activation function \cite{perezcorral2026lambdagelu}.\\

\noindent To strictly enforce this lower bound during unconstrained gradient-based optimization, the parameter $s \in \mathbb{R}$ is introduced and map it through the Softplus function:
\begin{equation*}
\lambda = 1 + \ln(1 + e^s)
\end{equation*}

\noindent Using this constrained parameter, the activation function is formulated by scaling the input within the cumulative distribution function:
\begin{equation*}
f(x) = x \Phi(\lambda x) = \dfrac{x}{2}\left[ 1 + \text{erf}\left(\dfrac{\lambda x}{\sqrt{2}}\right) \right]
\end{equation*}

\subsection{Gradient Formulation}
\noindent To ensure the proposed activation function is compatible with gradient-based optimization during backpropagation. To rigorously evaluate the derivative of the cumulative distribution function, we rely on the foundational principles of calculus.
\begin{theorem}[The First Fundamental Theorem of Calculus] 
If $f$ is continuous on $[a,b]$, then the function $g$ defined by
\begin{equation*}
g(x) = \int_{a}^{x}f(t)dt \quad a \leq x \leq b
\end{equation*}
is continuous on $[a,b]$ and differentiable on $(a,b)$, and $g'(x) = f(x)$.
\label{FFTC}
\end{theorem}
\noindent Extending Theorem~\ref{FFTC} with the Chain Rule, we have the following property: Suppose $f$ is continuous on $[a,b]$. If $v(x)$ is a differentiable function whose values lie within $[a,b]$ then,\\

\noindent If
$$
g(x) = \displaystyle\int_{a}^{v(x)} f(t)dt
$$
we have
$$
g'(x) = \dfrac{d}{dx}\left(  \displaystyle\int_{a}^{v(x)} f(t)dt\right) = f(v(x))\dfrac{d}{dx}(v(x))
$$
\noindent Given our formulation, we expand the error function into its integral form:
\begin{equation*}
f(x) = \frac{x}{2} \left[ 1 + \text{erf}\left( \frac{\lambda x}{\sqrt{2}}\right)\right] = \frac{x}{2} + \frac{x}{\sqrt{\pi}}\int_{0}^{\lambda x / \sqrt{2}} e^{-u^{2}}du
\end{equation*}
Observe that $e^{-u^{2}}$ is continuous on $\mathbb{R}$, and $\frac{\lambda x}{\sqrt{2}}$ is differentiable on $\mathbb{R}$. By direct application of the product rule and Theorem~\ref{FFTC}, we compute the derivative:
\begin{align*}
f'(x) &= \dfrac{d}{dx}\left( \frac{x}{2} + \frac{x}{\sqrt{\pi}}\int_{0}^{\lambda x / \sqrt{2}} e^{-u^{2}}du\right)\\ 
&= \dfrac{d}{dx}\left( \dfrac{x}{2}\right) + \dfrac{d}{dx} \left(  \frac{x}{\sqrt{\pi}}\int_{0}^{\lambda x / \sqrt{2}} e^{-u^{2}}du \right)\\
&=   \frac{1}{2} + \frac{1}{\sqrt{\pi}}\int_{0}^{\lambda x / \sqrt{2}} e^{-u^{2}}du + \frac{x}{\sqrt{\pi}}\frac{d}{dx}\int_{0}^{\lambda x / \sqrt{2}} e^{-u^{2}}du \\
&=   \frac{1}{2} + \frac{1}{\sqrt{\pi}}\int_{0}^{\lambda x / \sqrt{2}} e^{-u^{2}}du + \frac{x}{\sqrt{\pi}}e^{-(\lambda x)^{2} / 2}\left(\frac{\lambda}{\sqrt{2}}\right)
\end{align*}
\noindent which shows that the proposed activation function, $f(x)$, is continuously differentiable in $\mathbb{R}$.
\subsection{Asymptotic Behavior}
\noindent To demonstrate that the proposed activation function retains the core signal-preserving and noise-filtering properties of standard rectified units, we analyze its asymptotic behavior at the extremes.\\ 

\noindent First, we examine the behavior for large positive inputs. Given that our scaling parameter is strictly positive ($\lambda \geq 1$), as $x \to \infty$, the argument of the error function approaches $\infty$. By definition, the limit of the error function as its argument approaches $\infty$ is $1$. Evaluating the limit of our formulation:
\begin{align*}
\lim_{x \to \infty} f(x) &= \lim_{x \to \infty} \left( \frac{x}{2} \left[ 1 + \text{erf}\left( \frac{\lambda x}{\sqrt{2}} \right) \right] \right) \\
&= \left( \lim_{x \to \infty} \frac{x}{2} \right) \cdot \left( \lim_{x \to \infty} \left[ 1 + \frac{2}{\sqrt{\pi}} \int_{0}^{\lambda x / \sqrt{2}} e^{-u^2} du \right] \right) \\
&= \infty \cdot \left[ 1 + \frac{2}{\sqrt{\pi}} \int_{0}^{\infty} e^{-u^2} du \right] \\
&= \infty \cdot \left[ 1 + \frac{2}{\sqrt{\pi}} \left( \frac{\sqrt{\pi}}{2} \right) \right] \quad \left(\text{since } \int_{0}^{\infty} e^{-u^2} du = \frac{\sqrt{\pi}}{2}\right) \\
&= \infty \cdot (1 + 1) \\
&= \infty
\end{align*}
\noindent Because our function is defined as 
\begin{equation*}
f(x) = x\Phi(\lambda x)
\quad\Longrightarrow\quad
\frac{f(x)}{x} = \Phi(\lambda x),
\end{equation*}
where $x \neq 0$ for sufficiently large $x$ as $x \rightarrow +\infty$.
Since $\lambda \geq 1$, we have $\lambda x \rightarrow +\infty$, and hence
$\Phi(\lambda x) \rightarrow 1$. Therefore,
\begin{equation*}
\lim_{x\rightarrow +\infty}\frac{f(x)}{x}
=
\lim_{x\rightarrow +\infty}\Phi(\lambda x)
= 1.
\end{equation*}
Thus, $f(x) \sim x$ as $x \rightarrow +\infty$.\\

\noindent Although $\lim_{x\rightarrow+\infty}f(x)=+\infty$, the stronger result
$f(x)\sim x$ shows that the proposed activation is asymptotically linear
with unit slope on the positive domain. This is consistent with the
advantages of rectifier activations discussed by Glorot et al.
\cite{pmlr-v15-glorot11a}, who note that computation is linear on active
paths and that gradients flow well along these paths without the
activation-induced vanishing associated with sigmoid or $\tanh$ units.
Thus, the positive asymptotic behavior of the proposed activation
preserves an approximately identity-like mapping for large positive
inputs.\\

\noindent Next, we examine the behavior for large negative inputs. As $x \to -\infty$, given our scaling parameter $\lambda \geq 1$, the argument of the error function approaches $-\infty$. Evaluating the limit directly produces an indeterminate form:
\begin{equation*}
\lim_{x \to -\infty} \frac{x}{2} \left[ 1 + \text{erf}\left( \frac{\lambda x}{\sqrt{2}} \right) \right] \to -\infty \cdot 0
\end{equation*}
To resolve this, we rewrite the expression as a ratio and apply L'Hôpital's Rule~(Theorem~\ref{thm:lhospital}).
\begin{theorem}[L'H\^opital's Rule]
Let $f$ and $g$ be functions differentiable on an open interval $I$ containing $a$ except possibly at $a$ and $g'(x) \neq 0$ for all $x \in I \setminus \{a\}$.\\

\noindent If $\displaystyle\lim_{x \to a} \frac{f(x)}{g(x)}$ is indeterminate of type $\frac{0}{0}$ or $\frac{\infty}{\infty}$, then
$$
\lim_{x \to a} \frac{f(x)}{g(x)} = \lim_{x \to a} \frac{f'(x)}{g'(x)},
$$
provided $\displaystyle\lim_{x \to a} \frac{f'(x)}{g'(x)}$ exists, or is $+\infty$ or $-\infty$.
\label{thm:lhospital}
\end{theorem}
\noindent By Theorem~\ref{FFTC}  
\begin{align*}
\frac{d}{dz}\text{erf}(z) = \dfrac{d}{dz}\left(\dfrac{2}{\sqrt{\pi}}\int_{0}^{z} e^{-t^{2}}dt\right)= \frac{2}{\sqrt{\pi}}e^{-z^2}
\end{align*}
\noindent It follows that,
\begin{align*}
\lim_{x \to -\infty} f(x) &= \lim_{x \to -\infty} \frac{1 + \text{erf}\left( \frac{\lambda x}{\sqrt{2}} \right)}{2 x^{-1}} \\
&\overset{\text{L'HR}}{=} \lim_{x \to -\infty} \frac{ \frac{2}{\sqrt{\pi}} e^{-\frac{\lambda^2 x^2}{2}} \cdot \frac{\lambda}{\sqrt{2}} }{ -2 x^{-2} } \\
&= \lim_{x \to -\infty} \frac{ \frac{2}{\sqrt{2\pi}} \lambda e^{-\frac{\lambda^2 x^2}{2}} }{ -\frac{2}{x^2} } \\
&= \lim_{x \to -\infty} -\frac{\lambda}{\sqrt{2\pi}} \left( \frac{x^2}{e^{\frac{\lambda^2 x^2}{2}}} \right)
\end{align*}
As $x \to -\infty$, both $x^2$ and $e^{\frac{\lambda^2 x^2}{2}}$ approach $\infty$, yielding a $\frac{\infty}{\infty}$ indeterminate form. Applying L'Hôpital's Rule again:
\begin{align*}
\lim_{x \to -\infty} f(x) &\overset{\text{L'HR}}{=} \lim_{x \to -\infty} -\frac{\lambda}{\sqrt{2\pi}} \left( \frac{2x}{\lambda^2 x e^{\frac{\lambda^2 x^2}{2}}} \right) \\
&= \lim_{x \to -\infty} -\frac{\lambda}{\sqrt{2\pi}} \left( \frac{2}{\lambda^2 e^{\frac{\lambda^2 x^2}{2}}} \right) \\
&= 0
\end{align*}
\noindent Architecturally, the limit $\lim_{x \to -\infty} f(x) = 0$ acts as a soft-gating mechanism that induces network sparsity, similar to standard dropout. By suppressing negative inputs, it allows a Multi-Layer Perceptron (MLP) to efficiently isolate complex geometric features without activating the entire parameter space. Furthermore, because this collapse is asymptotic rather than instantaneous, it preserves a small gradient for negative values, mitigating the dead-neuron problem during optimization~\cite{pmlr-v15-glorot11a, hendrycks2023gaussianerrorlinearunits}.\\

\noindent Together, these limits mathematically guarantee that the parameterized activation effectively filters out significant negative noise while preserving strong positive signals, maintaining the fundamental advantages of ReLU while providing a continuously differentiable, smooth transition space around the origin.

\subsection{First Derivative Analysis and Gradient Stability}
\noindent To analyze the optimization dynamics and stability of our proposed activation function during backpropagation, we evaluate its first derivative obtained via differentiation rules:
\begin{equation*}
f'(x) = \Phi(\lambda x) + \lambda x \phi(\lambda x)
\end{equation*}

\noindent Let $f'(\cdot) = g(\cdot)$. Since $\lambda > 0$ by the Softplus parameterization ($\lambda = 1 + \ln(1 + e^s) \ge 1$), and the standard normal PDF is defined as $\phi(\lambda x) = \frac{1}{\sqrt{2\pi}} e^{-(\lambda x)^2/2}$.

\noindent We evaluate the asymptotic limits by distributing the limit operator:
\begin{equation*}
\lim_{x \rightarrow \pm \infty}\left(\Phi(\lambda x) + \lambda x\phi(\lambda x)\right) = \lim_{x \rightarrow \pm \infty} \Phi(\lambda x) + \lim_{x \rightarrow \pm \infty} \lambda x\phi(\lambda x)
\end{equation*} 

\noindent \textbf{As $x \rightarrow +\infty$:}
\begin{equation*}
\begin{aligned}
\lim_{x \rightarrow +\infty}\Phi(\lambda x) &= \lim_{x \rightarrow +\infty}\left(\dfrac{1}{2} + \dfrac{1}{\sqrt{\pi}}\int_{0}^{\lambda x/\sqrt{2}}e^{-u^{2}}du\right) = 1
\end{aligned}
\end{equation*}
Since $\displaystyle\lim_{x \rightarrow +\infty} \lambda x\phi(\lambda x) \rightarrow \frac{\infty}{\infty}$ is indeterminate, we apply Theorem~\ref{thm:lhospital}:
\begin{equation*}
\begin{aligned}
\lim_{x \rightarrow +\infty} \lambda x\phi(\lambda x) &= \lim_{x \rightarrow +\infty} \dfrac{\lambda}{\sqrt{2\pi}}\cdot\dfrac{x}{e^{(\lambda x)^2 / 2}}\\ 
&\overset{\text{L'HR}}{=} \lim_{x \rightarrow +\infty} \dfrac{\lambda}{\sqrt{2\pi}}\cdot\dfrac{1}{\lambda^2 x e^{(\lambda x)^2 / 2}} = 0
\end{aligned}
\end{equation*}
Thus, $\displaystyle\lim_{x \to +\infty} f'(x) = 1 + 0 = 1$.

\vspace{0.5cm}

\noindent \textbf{As $x \rightarrow -\infty$:}
\begin{equation*}
\begin{aligned}
\lim_{x \rightarrow -\infty}\Phi(\lambda x) &= \lim_{x \rightarrow -\infty}\left(\dfrac{1}{2} + \dfrac{1}{\sqrt{\pi}}\int_{0}^{\lambda x/\sqrt{2}}e^{-u^{2}}du\right) = 0
\end{aligned}
\end{equation*}
Similarly, $\displaystyle\lim_{x \rightarrow -\infty} \lambda x\phi(\lambda x) \rightarrow \frac{-\infty}{\infty}$ is indeterminate. Applying Theorem~\ref{thm:lhospital}:
\begin{equation*}
\begin{aligned}
\lim_{x \rightarrow -\infty} \lambda x\phi(\lambda x) &= \lim_{x \rightarrow -\infty} \dfrac{\lambda}{\sqrt{2\pi}}\cdot\dfrac{x}{e^{(\lambda x)^2 / 2}}\\ 
&\overset{\text{L'HR}}{=} \lim_{x \rightarrow -\infty} \dfrac{\lambda}{\sqrt{2\pi}}\cdot\dfrac{1}{\lambda^2 x e^{(\lambda x)^2 / 2}} = 0
\end{aligned}
\end{equation*}
Thus, $\displaystyle\lim_{x \to -\infty} f'(x) = 0 + 0 = 0$.

\begin{theorem}[Extreme Value Theorem]
\label{thm:evt}
If $f$ is continuous on a closed interval $[a,b]$, then $f$ attains an absolute maximum value $f(c)$ and an absolute minimum value $f(d)$ at some $c,d \in [a,b]$.
\end{theorem}

\noindent We need to show that $f'(x) = g(x) = \Phi(\lambda x) + \lambda x \phi(\lambda x)$ is globally bounded on $\mathbb{R}$, we leverage its continuity, asymptotic behavior, and Theorem~\ref{thm:evt}.\\

\noindent Because the standard normal CDF $\Phi(\cdot)$, the linear function $x$, and the standard normal PDF $\phi(\cdot)$ are continuous everywhere, their composition and sum $f'(x)$ is continuous on the entire real line $\mathbb{R}$.\\ 

\noindent From our previous asymptotic analysis, we established that:
\begin{equation*}
\lim_{x \to -\infty} f'(x) = 0 \quad \text{and} \quad \lim_{x \to +\infty} f'(x) = 1
\end{equation*}
By the definition of limits at infinity, there exists a sufficiently large real number $M > 0$ such that for all $x > M$, $f'(x)$ is bounded arbitrarily close to $1$, and for all $x < -M$, $f'(x)$ is bounded arbitrarily close to $0$. Thus, on the unbounded open intervals $(-\infty, -M)$ and $(M, +\infty)$, $f'(x)$ is bounded.\\

\noindent Within the finite closed interval $[-M, M]$, $f'(x)$ is continuous. By the Extreme Value Theorem (Theorem~\ref{thm:evt}), $f'(x)$ must attain an absolute maximum and minimum within $[-M, M]$, guaranteeing it is bounded on this closed interval.\\

\noindent Since $f'(x)$ is bounded on the closed interval $[-M, M]$ and is constrained by finite limits on the external regions $(-\infty, -M)$ and $(M, +\infty)$, it follows that $f'(x)$ is strictly bounded for all $x \in \mathbb{R}$. Because
$$
\displaystyle\lim_{x \to -\infty} f'(x) = 0 \text{ and } \displaystyle\lim_{x \to +\infty} f'(x) = 1$$
the function cannot escape to infinity in either direction. Thus, there exist a constant $K > 0$ such that:
\begin{equation*}
\lvert f'(x) \rvert \le K \quad \forall x \in \mathbb{R}
\end{equation*}
This bounded derivative guarantees that MPGELU maintains gradient stability, mathematically ensuring that error signals do not compound exponentially and preventing the exploding gradient problem during backpropagation.\\

\vspace{0.5cm}

\noindent While the bounded derivative $|f'(x)| \le K$ prevents gradient explosion, we must also ensure that gradients do not vanish. Vanishing gradients typically occur under two conditions: saturation at large positive values (as seen in Sigmoid or Tanh) and dead zones near the origin (as seen in standard ReLU).\\

\noindent Modern neural networks initialize weights using zero-mean distributions, concentrating early pre-activation values near the origin ($x \approx 0$). In standard ReLU, inputs slightly below zero receive a gradient of exactly $0$, risking permanent deactivation. Evaluating the MPGELU derivative at the origin yields:
\begin{equation*}
f'(0) = \Phi(0) + 0 \cdot \phi(0) = 0.5
\end{equation*}
\noindent Because $f'(0) > 0$, MP-GELU guarantees active gradient flow right from initialization. This non-zero derivative at the origin effectively avoids the "dying neuron" pathology, ensuring continuous and stable learning updates.



\bibliography{references.bib}
\end{document}